% ============================================================
% preamble.tex —— 大连理工大学本科毕业论文 LaTeX 格式设置
% 编译方式：xelatex -> biber -> xelatex -> xelatex
% ============================================================

% ---------- 页面尺寸 ----------
\usepackage{geometry}
\geometry{
  paper     = a4paper,
  top       = 35mm,
  bottom    = 25mm,
  left      = 25mm,
  right     = 25mm,
  headsep   = 10mm,
  footskip  = 20mm,
  headheight= 14pt,
}

% ---------- 中英文字体 ----------
\usepackage{fontspec}
\usepackage{xeCJK}
% 中文字体（Windows系统内置）
\setCJKmainfont{SimSun}[
  BoldFont   = SimHei,
  ItalicFont = KaiTi,
  AutoFakeBold = 2,
]
\setCJKsansfont{SimHei}[AutoFakeBold = 2]
\setCJKmonofont{FangSong}
% 华文细黑（封面题目用，AutoFakeBold模拟加粗）
\newfontfamily\xihei{STXihei}[AutoFakeBold=2]
% 英文字体
\setmainfont{Times New Roman}
\setsansfont{Arial}
\setmonofont{Courier New}

% ---------- 字号别名（方便正文调用）----------
% \zihao{-4} = 小四号 = 12pt（正文默认）
% \zihao{4}  = 四号   = 14pt（节标题）
% \zihao{-3} = 小三号 = 15pt（章标题）

% ---------- 行距 ----------
\usepackage{setspace}
\setstretch{1.25}                     % 多倍行距1.25（模板要求）

% ---------- 首行缩进 ----------
\usepackage{indentfirst}
\setlength{\parindent}{2\ccwd}        % 首行缩进2个中文字符

% ---------- 段间距 ----------
\setlength{\parskip}{0pt}

% ---------- 章节标题格式（ctex宏包已在主文件加载）----------
\ctexset{
  % 一级标题：黑体，小三，居左，段后1行
  section = {
    format      = \zihao{-3}\heiti\raggedright,
    beforeskip  = 0pt,
    afterskip   = 22pt,   % 1行 = 小三(15pt)×1.5倍行距 ≈ 22pt
  },
  % 二级标题：黑体，四号，居左
  subsection = {
    format      = \zihao{4}\heiti\raggedright,
    beforeskip  = 6pt,
    afterskip   = 0pt,
  },
  % 三级标题：黑体，小四，居左
  subsubsection = {
    format      = \zihao{-4}\heiti\raggedright,
    beforeskip  = 6pt,
    afterskip   = 0pt,
  },
}

% ---------- 目录格式 ----------
\usepackage{tocloft}
% 目录标题：黑体，小三，居中
\renewcommand{\contentsname}{目\quad 录}
% 一级目录：黑体，小四
\renewcommand{\cftsecfont}{\heiti\zihao{-4}}
\renewcommand{\cftsecpagefont}{\zihao{-4}}
% 二级及以下目录：宋体，小四
\renewcommand{\cftsubsecfont}{\songti\zihao{-4}}
\renewcommand{\cftsubsecpagefont}{\zihao{-4}}
\renewcommand{\cftsubsubsecfont}{\songti\zihao{-4}}
\renewcommand{\cftsubsubsecpagefont}{\zihao{-4}}
\setlength{\cftbeforesecskip}{6pt}

% ---------- 页眉页脚 ----------
\usepackage{fancyhdr}
\pagestyle{fancy}
\fancyhf{}
\fancyhead[C]{\zihao{5}\songti 基于在线评论的消费者多属性态度预测效度研究}
\fancyfoot[C]{\zihao{-5}\songti \thepage}
\renewcommand{\headrulewidth}{0.4pt}
\renewcommand{\footrulewidth}{0pt}

% 纯页码样式（摘要、目录用）
\fancypagestyle{plain}{
  \fancyhf{}
  \fancyfoot[C]{\zihao{-5}\songti \thepage}
  \renewcommand{\headrulewidth}{0pt}
}

% ---------- 参考文献（GB/T 7714-2015）----------
\usepackage[
  backend  = biber,
  style    = gb7714-2015,
  citestyle = gb7714-2015,
  gbpub    = false,
  doi      = false,
  url      = false,
]{biblatex}
\addbibresource{refs/refs01.bib}   % [1]-[10]  态度理论、评论影响
\addbibresource{refs/refs02.bib}   % [11]-[20] 文本信息量、行为意向
\addbibresource{refs/refs03.bib}   % [21]-[30] CMV/效度、ABSA、LLM
\addbibresource{refs/refs04.bib}   % [31]-[41] 中文NLP、消费者专业性
\addbibresource{refs/refs05.bib}   % [42]+  补充文献条目

% ---------- 表格 ----------
\usepackage{booktabs}     % 三线表：\toprule \midrule \bottomrule
\usepackage{longtable}    % 跨页长表格
\usepackage{array}
\usepackage{multirow}
\usepackage{tabularx}
\usepackage{makecell}
\usepackage{threeparttable}  % tablenotes 表格注脚环境

% 表格字体默认小五号
\let\oldtabular\tabular
\renewcommand{\tabular}{\zihao{5}\oldtabular}

% ---------- 图片 ----------
\usepackage{graphicx}
\usepackage{float}
\graphicspath{{figures/}}

% ---------- 图表标题格式 ----------
\usepackage{caption}
\captionsetup{
  font       = {small},             % 宋体五号
  labelsep   = space,
  skip       = 6pt,
}
% 图标题在图下方，表标题在表上方（通过放置位置控制，无需captionsetup）

% 图/表编号格式：图1-1，表2-3
\renewcommand{\thefigure}{\thesection-\arabic{figure}}
\renewcommand{\thetable}{\thesection-\arabic{table}}

% ---------- 数学公式 ----------
\usepackage{amsmath}
\usepackage{amssymb}
\usepackage{amsthm}

% 公式编号格式：(1-1)
\renewcommand{\theequation}{\thesection-\arabic{equation}}

% ---------- tikz（研究框架图）----------
\usepackage{tikz}
\usetikzlibrary{
  shapes.geometric,
  arrows.meta,
  positioning,
  calc,
  fit,
  backgrounds,
  decorations.pathreplacing,
}

% ---------- 超链接（目录可点击）----------
\usepackage[
  colorlinks = true,
  linkcolor  = black,
  citecolor  = black,
  urlcolor   = blue,
  hidelinks,
]{hyperref}

% ---------- 其他常用宏包 ----------
\usepackage{xcolor}
\usepackage{listings}       % 代码块（附录用）
\usepackage{enumitem}       % 列表格式控制
\usepackage{footnotehyper}  % 脚注
\usepackage{rotating}       % 旋转表格
\usepackage{pdflscape}      % 横向页面

% 列表间距压缩
\setlist{
  topsep    = 4pt,
  itemsep   = 2pt,
  parsep    = 0pt,
  leftmargin = 2\ccwd,
}

% ---------- 代码块样式 ----------
\lstset{
  basicstyle   = \ttfamily\small,
  breaklines   = true,
  frame        = single,
  numbers      = left,
  numberstyle  = \tiny\color{gray},
  keywordstyle = \color{blue}\bfseries,
  commentstyle = \color{gray}\itshape,
  stringstyle  = \color{orange},
  showstringspaces = false,
  xleftmargin  = 2em,
}

% ---------- 自定义命令 ----------
% 统计符号
\newcommand{\N}{\textit{N}}
\newcommand{\pval}{\textit{p}}
\newcommand{\betacoef}{\(\beta\)}
\newcommand{\Rsq}{\(R^2\)}
\newcommand{\DeltaRsq}{\(\Delta R^2\)}

% ATT公式
\newcommand{\ATTformula}{%
  \text{ATT} = \sum_{i=1}^{n} b_i \times e_i
}

% 显著性标注
\newcommand{\sig}[1]{%
  \ifcase#1
    \or $^{*}$%
    \or $^{**}$%
    \or $^{***}$%
  \fi
}

% 假设环境
\newcounter{hypothesis}
\newenvironment{hypothesis}[1][]{%
  \refstepcounter{hypothesis}%
  \par\medskip\noindent%
  \textbf{假设\thehypothesis（H\thehypothesis）}%
  \ifx&#1&\else（#1）\fi：%
}{%
  \par\medskip%
}

% 注释行（表格下方）
\newcommand{\tablenote}[1]{%
  \par\smallskip\noindent\footnotesize\textit{注：}#1
}

% ---------- PDF元数据 ----------
\hypersetup{
  pdftitle  = {基于在线评论的消费者多属性态度预测效度研究},
  pdfauthor = {大连理工大学},
}
